\section{Conclusion}\label{sec:conclusion}

Ce papier demandait si une baisse de la fiscalité des dividendes fait entrer du capital dans les
entreprises ou nourrit seulement la distribution aux ménages, et si la réponse du capital à l'impôt
est symétrique. La réforme française de 2018, le prélèvement forfaitaire unique, et la hausse
d'impôt de 2013 qui l'avait précédée sur la même marge fournissent le double choc qui permet d'y
répondre.

Trois résultats ressortent, chacun répondant à une hypothèse. Sur la première, la baisse d'impôt de
2018 ne fait entrer aucun capital neuf : le coefficient sur la variation du capital souscrit,
$-0{,}0034$, est un nul estimé avec précision, et la permutation placebo place la France au plus
faible écart des six pays. Sur la deuxième, aucune substitution par la dette ni effet réel
n'apparaît ; la complémentarité comptable entre part de fonds propres et part de dettes rend les
deux coefficients égaux au signe près et tous deux nuls, de sorte qu'au point de fonds propres perdu
ne répond aucun point de dette gagné. Ces deux résultats confirment la neutralité des nouvelles vues
de l'impôt sur les dividendes, dans laquelle la fiscalité de la distribution ne déplace pas la marge
d'investissement. L'argent libéré par la réforme sort vers les ménages, dont les dividendes reçus
bondissent de 32 \% l'année même, et se loge en épargne liquide et en actions de portefeuille, non
dans le capital des sociétés privées ni dans l'immobilier. Cet écart entre la vive réponse de la
distribution et l'absence de réponse réelle rejoint ce que \citet{chetty2005} et \citet{yagan2015}
observent sur la baisse américaine de 2003.

Le résultat central est le troisième. En plaçant la hausse de 2013 et la baisse de 2018 dans un
cadre d'estimation unique, nous mettons au jour un cliquet : sur la distribution, la baisse ne défait
que 45 \% de l'effet de la hausse ; sur l'allocation du capital des sociétés privées, elle n'en
ramène qu'environ un tiers, la réallocation revenant dans un rapport de un à près de trois sur les
réserves comme sur les fonds propres. Le capital part plus facilement qu'il ne revient. La
contribution du papier tient à ce contraste placé dans un cadre unique, qui fait de l'asymétrie
elle-même l'objet de mesure là où la littérature n'observait qu'un choc isolé, et à sa cohérence sur
deux niveaux indépendants, la distribution agrégée des ménages et l'allocation du capital des
sociétés.

L'implication pour l'évaluation de la fiscalité du capital est directe et forte. Une baisse d'impôt
symétrique à une hausse passée ne défait pas ses effets : les déformations qu'une hausse imprime à
l'allocation du capital persistent après la baisse qui lui succède, parce que les coûts d'ajustement
et les motifs de rétention rendent la sortie des fonds plus réactive que leur retour. Un calcul
d'évaluation qui suppose la réversibilité surestime donc le gain d'une baisse et sous-estime le coût
durable d'une hausse.

La portée du résultat est bornée par le cadre d'identification. L'inférence repose sur six pays, ce
qui limite la puissance de la permutation placebo, et la trajectoire du capital retenu est mesurée
relativement aux contrôles, de sorte qu'elle établit que l'élévation postérieure à 2013 est propre à
la France sans l'imputer à la seule réforme de 2013 à l'exclusion de tout autre fait français de la
même année. Ces limites invitent à répliquer le contraste des deux chocs sur d'autres pays ayant
parcouru la même marge fiscale dans les deux sens, que les données microéconomiques de sociétés
privées rendent désormais accessibles.
